\scanpage{27}
\mspage{25}
% archive sheet number: 26
Mais ce morphisme est-il un \uline{iso}morphisme ? Je pense que dans le cas additif ça peut être raisonnable, mais dans le cas non additif ça n'a aucune raison d'être vrai. On voit en tous cas que si pour tout $f:I\longrightarrow I'$ dans $\mathrm{Diag}$, $f_*^{\mathbb D}:\mathbb D_I\longrightarrow\mathbb D_{I'}$ est (non seulement exact à g., mais aussi) exact à droite, alors pour tout $g:J\longrightarrow J'$ dans $\mathrm{Diag}$, $g_!^{\mathbb D}$ est (non seulement exact à droite, mais même) exact à gauche.

Prenons le dérivateur des préfaisceaux d'ens. sur des $I$ variables, sa cat. fondamentale est $\mathrm{Ens}$. \inserted{c'est un dérivateur non pointé} Considérons $I\longrightarrow I'$ une imm. fermée \deleted{\unsure{et}} \inserted{non \deleted{\illegible{1}} surjective} $f_{\unsure{*|!}}:\mathbb D(I)\longrightarrow\mathbb D(I')$. Si $f_*^{\mathbb D}$ était \unsure{exacte|exact} à droite, \deleted{il} \inserted{$f_*$} devrait transformer objet initial de $\mathbb D(I)$ en objet initial de $\mathbb D(\unsure{I'|I})$, sauf erreur. (Cf plus bas) Or il n'en est rien, puisque pour \unsure{tout} $F\in\mathrm{Im}(f_*^{\mathbb D}:\mathbb D(I)\longrightarrow\mathbb D(I'))$ la restriction de $F$ à $I'\setminus I=U$ est le faisceau final $e_U=(i\longmapsto\{e\})$, et non le faisceau initial $\varnothing_U=(i\longmapsto\varnothing)$.

\unsure{Donc on prouvera ceci|On va prouver ceci} :

\uline{Prop} : \deleted{\unsure{Si un}} $\mathbb D$ est un dérivateur à g. (à dr.), alors les produits finis (sommes finies) existent dans les $\mathbb D(I)$. Tout \deleted{\unsure{f}} morphisme \deleted{de} $\varphi:\mathbb D\unsure{\longrightarrow|-}\mathbb D'$ de dérivateurs qui est exact à g (à dr.) commute auxdits produits (sommes), \deleted{\unsure{ceci}} pour chaque $I$ donnant $\varphi(I):\mathbb D(I)\longrightarrow\mathbb D'(I)$. En particulier il transforme objet final (initial) en \unsure{itou|idem}
